\chapter{Introduction}
The adoption of Electronic Health Records (EHRs) has grown exponentially over the past decade, resulting in a massive digitization of patient data across global healthcare networks. While structured data points like vital signs are easily managed, the vast majority of clinical information is captured as unstructured text in the form of medical transcriptions, physician notes, and discharge summaries. The increasing volume of this unstructured clinical text presents a critical bottleneck for hospital administrators who must efficiently process, categorize, and route these documents to specialized departments. Natural Language Processing (NLP) has emerged as an indispensable technology in this domain, providing the computational tools necessary to automate the analysis of complex medical narratives. By bridging the gap between human language and machine understanding, NLP enables the automated classification of these documents, naturally transitioning medical informatics from slow, manual sorting to intelligent, algorithm-driven routing systems.

\section{Background}
The rapid digitization of global healthcare systems has resulted in the unprecedented accumulation of clinical data. Electronic Health Records (EHRs) and medical transcriptions encompass comprehensive patient histories, diagnoses, and treatment plans. However, a significant portion of this invaluable clinical documentation is stored in unstructured, narrative text formats. Medical transcription, the process of converting voice-recorded medical reports dictated by healthcare practitioners into text, generates vast quantities of data. Traditionally, these transcription reports are manually reviewed and organized by specific medical specialties, such as Cardiology, Neurology, or Orthopedics. 

\section{Motivation}
The primary motivation for this project is to address the severe inefficiencies associated with the manual handling of healthcare documents. As patient populations and data volumes grow exponentially, manual classification methods become unsustainable. Automating this process using Natural Language Processing (NLP) and Machine Learning (ML) can effectively reduce administrative overhead, ensure consistent document routing, and allow healthcare professionals to redirect their focus from clerical tasks directly to patient care. To improve classification robustness and reduce the limitations of relying on a single classifier, this project proposes a Voting Ensemble approach that combines multiple machine learning models.

\section{Problem Statement}
In contemporary healthcare institutions, organizing unstructured medical documents remains a substantial administrative challenge. Relying on manual categorization by hospital personnel is inherently slow, financially expensive, and subject to natural human error. The subjective interpretation of complex medical jargon frequently leads to inconsistent categorizations, which can misroute critical patient information and severely delay clinical workflows. Therefore, the core problem is the lack of an efficient, scalable, and highly accurate automated mechanism capable of classifying unstructured clinical text into precise medical specialties.
